\documentclass[10pt]{article}
\usepackage[utf8]{inputenc}
\usepackage[T1]{fontenc}
\usepackage{amsmath}
\usepackage{amsfonts}
\usepackage{amssymb}
\usepackage[version=4]{mhchem}
\usepackage{stmaryrd}
\usepackage{bbold}
\usepackage{graphicx}
\usepackage[export]{adjustbox}
\graphicspath{ {./images/} }
\usepackage{hyperref}
\hypersetup{colorlinks=true, linkcolor=blue, filecolor=magenta, urlcolor=cyan,}
\urlstyle{same}

\begin{document}
Moral hazard-Principal agent model

\section*{Principal-Agent. Introduction}
\begin{itemize}
  \item The principal ˋ"controls" an operation in which the agent participates.
  \item The agent's actions have consequences for the principal.
  \item The principal does not observe perfectly the agent's actions.
  \item The principal designs an incentive system to best align the agent's actions with the principal's objectives.
  \item Example: An owner (principal) hires an agent to manage a restaurant; the manager may work hard or not; the owner cannot perfectly monitor the manager's effort.
  \item Moral hazard: Asymmetry of information arises from a "hidden" actionan agent's action that is not perfectly observable by the principal.
\end{itemize}

Adverse selection: Asymmetry of information arises from information only available to one party-private information.

\section*{Applications}
\begin{itemize}
  \item See Bolton and Dewatripont (2005).
  \item Survey, The theory of contracts, Hart and Holmström (1987).
  \item Managerial theory of the firm; Jensent and Meckling (1976), Grossman and Hart (1983).
  \item Optimal contracts between landlords and tenants, sharecroping.
  \item Efficiency wage theories.
  \item Financial intermediation; Diamond (1984).
  \item Theoretical Accounting; Demski and Kreps (1982).
\end{itemize}

\section*{Models}
There are at least three representations of the agency model (Hart and Holmström (1987)).

\begin{itemize}
  \item The state space formulation,
  \item The parametrized formulation,
  \item The abstract formulation.
\end{itemize}

\section*{State space formulation, a worker}
\begin{itemize}
  \item The owner of a firm, the principal, hires a worker or manager, the agent.
  \item The firm's technology produces a publicly observable output $y$.
  \item $y=y(\theta, e)$ where $y$ : $\Theta \times A \rightarrow \mathbb{R}$.
\end{itemize}

For simplicity $y=\theta+e$.

\begin{itemize}
  \item $\theta \in \mathbb{R}$, unobservable random variable with density $f(\theta)>0$
  \item $e \in A \subset \mathbb{R}$ agent's effort level, unobservable to principal.
  \item $t: \mathbb{R} \rightarrow \mathbb{R}$, money transfer $t(y)$ from the principal to the agent when output is $y$; compensation package.
  \item Principal's utility: $U_{p}(y-t(y)), U_{p}^{\prime}>0, U_{p}^{\prime \prime} \leq 0$.
  \item Agent's utility: $U_{a}(t)-\Psi(e)$, where $U_{a}^{\prime}>0, U_{a}^{\prime \prime} \leq 0$ (concave), and $\Psi^{\prime}>0$ and $\Psi^{\prime \prime} \geq 0$ (convex).\\
$\Psi(e)$ is the cost to the agent of exerting effort $e$.\\
It is common in the literature to assume that the agent's utility is separable in income and effort.
  \item $\bar{U}_{a}$ is the agent's reservation utility.
\end{itemize}

\section*{The principal-agent model-Timing}
\begin{enumerate}
  \item The principal proposes a contract $t(y)$ where $y$ is observable output and $t$ is the payment to the agent.
  \item If the agent rejects the contract, the game ends. Players receive their reservation utility.
  \item If agent accepts, agent chooses his action (effort level) $e \in A$. The choice $e$ is not contractable.
  \item Output $y(\theta, e)$ is realized and observed, where $\theta$ is an unobserved random variable with known distribution.
  \item The agent receives $t(y)$, the principal receives $y-t(y)$ and the game ends.
\end{enumerate}

\section*{Complete information. Benchmark}
\begin{itemize}
  \item The principal's problem is to find a function $t$ (that depends only on the observable variable $y$ and that maximizes the principal's utility inducing the agent to participate in the process (i.e. IR).
  \item For comparison, it is useful to consider the complete information case, i.e. the effort level $e \in A$ is observable and contractable.
\end{itemize}

\section*{Principal's problem, full information}
$$
\begin{aligned}
\max _{t, e} & E_{\theta} U_{p}(y-t(y)) \\
\text { s.t. } & y=y(\theta+e) \\
(I R) \quad & E\left[U_{a}(t)-\Psi(e)\right] \geq \bar{U}_{a} \\
& e \in A \\
& t \in L^{1}(\mathbb{R})
\end{aligned}
$$

In an equivalent formulation the Principal chooses a transfer function $t(y, e)$.

For definition of $L_{1}$, see for instance Royden (1968).

\section*{First order conditions}
$$
\begin{aligned}
\mathcal{L}(t, e, \lambda)= & \int U_{p}(y-t(y)) f(\theta) \mathrm{d} \theta \\
& +\lambda\left[\int U_{a}(t(y)) f(\theta) \mathrm{d} \theta-\Psi(e)-\bar{U}_{a}\right]
\end{aligned}
$$

First order conditions with respect to $t$ and $e$,


\begin{gather*}
0=-U_{p}^{\prime} f+\lambda U_{a}^{\prime} f  \tag{1}\\
0=\int U_{p}^{\prime} \times \underbrace{\left(\frac{\partial y(\theta, e)}{\partial e}\right)}_{\frac{\partial(\theta+e)}{\partial e}=1} f \mathrm{~d} \theta-\lambda \Psi^{\prime}(e) \tag{2}
\end{gather*}


FOC with respect to $t$ and $e$,

$$
\begin{aligned}
0 & =-U_{p}^{\prime} f+\lambda U_{a}^{\prime} f \\
U_{p}^{\prime} & =\lambda U_{a}^{\prime}
\end{aligned}
$$

where $f$ is eliminated because $f>0$.

$$
0=\int U_{p}^{\prime} f \mathrm{~d} \theta-\lambda \Psi^{\prime}(e)
$$

\section*{Aside: order of derivation and integration}
Let $f: X \times A \rightarrow \mathbb{R}$ such that $\forall x \in X, f(x, \cdot)$ is differentiable with respect to $a$; and $\forall a \in A, f(\cdot, a)$ is integrable. For a distribution $\mu$ on $X$, define $G(a)=\int f(x, a) \mathrm{d} \mu$, then $\frac{\mathrm{d} G}{\mathrm{~d} a}=\lim _{h \rightarrow 0} \frac{G(a+h)-G(a)}{|h|}$.

$$
\begin{aligned}
\frac{\mathrm{d} G}{\mathrm{~d} a} & =\lim _{h \rightarrow 0} \int \frac{f(x, a+h)-f(x, a)}{|h|} \mathrm{d} \mu \\
& =? \int \lim _{h \rightarrow 0} \frac{f(x, a+h)-f(x, a)}{|h|} \mathrm{d} \mu \\
& =\int \frac{\mathrm{d} f(x, a)}{\mathrm{d} a} \mathrm{~d} \mu
\end{aligned}
$$

\begin{itemize}
  \item Let $g^{h}(x, a)=\frac{f(x, a+h)-f(x, a)}{|h|}$.
  \item Since $f$ is differentiable wrt $a$, for fix $a \lim _{h \rightarrow 0} g^{h}(x, a)=\frac{\mathrm{d} f(x, a)}{\mathrm{d} a}$ pointwise.
  \item The Dominated Convergence Theorem completes the argument.
\end{itemize}

\section*{Aside: Gateaux directional derivative}
Definition $\mathbf{1}$ (Gateaux differential) Let $X$ be a vector space, $Y$ a normed space, and $T: X \rightarrow Y$. The Gateaux differential of $T$ at $x \in X$ with increment $h \in X$ is

$$
\gamma T(x ; h)=\lim _{\alpha \downarrow 0} \frac{1}{\alpha}[T(x+\alpha h)-T(x)]
$$

if it exists.\\
If the limit exists for every $h \in X$, the function $T$ is said to be Gateaux differentiable at $x$.

\section*{RN principal, RA agent, observable $e$}
Risk-neutral principal, risk-averse agent, observable $e$.

\begin{itemize}
  \item Wlog let $U_{p}(y-t)=y-t(y)$ and $U_{p}^{\prime}=1$.
\end{itemize}

Then Equation 1, FOC w.r.t $t$, becomes

$$
\begin{aligned}
-U_{p}^{\prime}+\lambda U_{a}^{\prime} & =0 \\
\lambda & =\frac{U_{p}^{\prime}}{U_{a}^{\prime}} \\
U_{a}^{\prime}(t(y)) & =\frac{1}{\lambda} \forall y
\end{aligned}
$$

\begin{itemize}
  \item Then optimal contract $t(y)=k$, a constant.
\end{itemize}

FOC w.r.t. $e$ used to determine the optimal $e^{*}$.

$$
\begin{aligned}
0 & =\int U_{p}^{\prime} \times\left(\frac{\partial y(\theta, e)}{\partial e}\right) f \mathrm{~d} \theta-\lambda \Psi^{\prime}(e) \\
0 & =1-\lambda \Psi^{\prime}(e) \text { because } U_{p}^{\prime}=1 \\
\Psi^{\prime}(e) & =\frac{1}{\lambda}
\end{aligned}
$$

Since $\frac{1}{\lambda}=U_{a}^{\prime}(t(y))$, the optimal effort level $e^{*}$ is the solution to

$$
\Psi^{\prime}\left(e^{*}\right)=U_{a}^{\prime}(t(y))=U_{a}^{\prime}(k)
$$

\section*{Inefficiency of $t(y)=k$ with unobservable $e$}
$$
\begin{aligned}
& \max _{e \in A} E U_{a}(t)-\Psi(e) \\
& \max _{e \in A} E U_{a}(k)-\Psi(e)
\end{aligned}
$$

The FOC wrt to " $e$ " is

$$
-\Psi^{\prime}(e)=0
$$

QED

\begin{itemize}
  \item The optimal contract under imperfect information was not derived.
  \item Thus, previous slide does not "prove" that the first best cannot be achieved under imperfect information.
  \item That conclusion follows, however, because the only way to provide complete insurance to the risk-averse agent is with a constant $t(y)=k$.
  \item Constant transfer would necessarily elicit lower effort level than it is optimal under complete information.
\end{itemize}

\section*{RA principal, RN agent, observable $e$}
\begin{itemize}
  \item W.I.o.g. let $U_{a}(t(y))=t(y)$ and $U_{a}^{\prime}=1$.
\end{itemize}

From FOC Equation 1 , since $f>0$, it follows that

$$
\forall y, U_{p}^{\prime}(y-t(y))=\lambda
$$

This implies that $y-t(y)=K$, a constant.

\begin{itemize}
  \item From FOC Equation 2, $1-\Psi^{\prime}\left(e^{*}\right)=0$.
  \item Single instrument and single objective, risk sharing. With incomplete information there are two objectives, risk sharing and correct incentives to work.
\end{itemize}

Theorem 1 Suppose the effort level $e \in A$ is not observable by the principal, the principal is risk averse (RA-P) and the agent is risk neutral (RN-A). Then, there exists a contract that achieves the first best solution, i.e., the solution of the complete information case.

\begin{itemize}
  \item Intuition: Since the agent is risk neutral, risk sharing is not an issue. Thus $t(y)$ need only elicit the correct effort level from the agent.
\end{itemize}

\section*{RA-Principal, RN-A, unobservable $e$}
\begin{itemize}
  \item Let $e^{*}, t^{*}(y)$ be the optimal solution under complete information.
  \item Then, $y-t^{*}(y)=K$ and therefore $t^{*}(y)=y-K$.
  \item The value of $K$ must be determined.
\end{itemize}

The value of $K$ is obtained from IR.

$$
\begin{aligned}
E U_{a}(t(y))-\Psi\left(e^{*}\right) & =\bar{U}_{a} \\
E U_{a}(t(y)) & =\Psi\left(e^{*}\right)+\bar{U}_{a} \\
E t(y) & =\Psi\left(e^{*}\right)+\bar{U}_{a} \\
E y-K & =\Psi\left(e^{*}\right)+\bar{U}_{a} \\
e^{*}+E \theta-K & =\Psi\left(e^{*}\right)+\bar{U}_{a} \\
e^{*}+E \theta-\Psi\left(e^{*}\right)-\bar{U}_{a} & =K
\end{aligned}
$$

Therefore

$$
t^{*}(y)=y-K=y+\bar{U}_{a}+\Psi\left(e^{*}\right)-e^{*}-E \theta
$$

It remains to verify that with $t^{*}$, agent chooses $e^{*}$.

$$
\begin{aligned}
& \max _{e \in A} E U_{a}(t)-\Psi(e) \\
& \max _{e \in A} E t^{*}(y)-\Psi(e) \\
& \max _{e \in A} E y+\bar{U}_{a}+\Psi\left(e^{*}\right)-e^{*}-E \theta-\Psi(e) \\
& \max _{e \in A} E \theta+e+\bar{U}_{a}+\Psi\left(e^{*}\right)-e^{*}-E \theta-\Psi(e)
\end{aligned}
$$

The FOC wrt to " $e$ " is, $1-\Psi^{\prime}(e)=0$.\\
QED

\section*{Preliminaries.}
\begin{itemize}
  \item For the definitions and basic results on stochastic dominance see, for instance, Mas-Colell et al. (1995), page 195-197.
  \item First- and second-order stochastic dominance were introduced to economics in Rothschild and Stiglitz (1976).
  \item See also Marshall et al. (2010), Chapter 17, pages 707-712.
\end{itemize}

\section*{FOSD}
Definition 2 (First-order stochastic dominance) Let $\mu, \nu$ be two probability measures on $X \subset \mathbb{R}$.

We say that $\mu$ first-order stochastically dominates $\nu$ if for every nondecreasing function $u: \mathbb{R} \rightarrow \mathbb{R}, \int u(x) d \mu \geq \int u(x) d \nu$.

Theorem 2 The cummulative distribution function (cdf) $F$ first-order stochastically dominates the cdf $G$ if and only if $(\forall x) F(x) \leq G(x)$.

\section*{FOSD}
$\frac{1}{\frac{x}{1}}$\\
$\frac{1}{\frac{x}{10}}$\\
1.0\\
0.9\\
0.8\\
0.7\\
0.6\\
0.5\\
0.5\\
0.4\\
0.4\\
0.3\\
0.3\\
"

$$
\begin{aligned}
& F>_{F O S D} G \\
& (\forall x), 1-F(x)>1-G(x)
\end{aligned}
$$

\section*{SOSD}
Definition 3 (Second-Order Stochastic Dominance) Let $F, G$ be two cdf with the same mean. $F$ second-order stochastically dominates (or is less risky than) $G$ if for every nondecreasing concave function $u: \mathbb{R}_{+} \rightarrow \mathbb{R}$, $\int u(x) d F(x) \geq \int u(x) d G(x)$.

Definition 4 (Mean-Preserving Spread) Let $F$ be the cdf corresponding to a random variable $x$. For any realization $x$, let $H_{x}(z)$ be the cdf corrsponding to a random variable $z$ and suppose $\int z d H_{x}(z)=0$. Consider the new random variable $y=x+z(x)$ and let $G$ be its cdf. Then we say that $G$ is a mean-preserving spread of $F$.

Theorem 3 Let $F, G$ be two cdfs with the same mean. The following statements are equivalent.

\begin{enumerate}
  \item $F$ second-order stochastically dominates $G$.
  \item $G$ is a mean-preserving spread of $F$.
  \item $\int_{0}^{x} G(t) d t \geq \int_{0}^{x} F(t) d t, \forall x$.
\end{enumerate}

\section*{Parametrized distribution approach}
\begin{itemize}
  \item It provides intuitive necessary conditions for the optimal mechanism.
  \item Output is a random variable $y \in[\underline{y}, \bar{y}]$ with twice continuously differentiable cdf $F(y, e)$ and density $f(y, e)$.
  \item First order stochastic dominance FOSD implies that $e^{\prime}>e \Longrightarrow F\left(x, e^{\prime}\right) \leq F(x, e)$.
\end{itemize}

Higher values of $e$ yield probabilistically higher $y$.

\begin{itemize}
  \item $\frac{f\left(y, e^{\prime}\right)}{f(y, e)}$ must be bounded, see Mirrlees (1999).
\end{itemize}

\section*{FOSD in parametrized model}
\includegraphics[max width=\textwidth, alt={}, center]{3d1c6c7a-b07a-4a39-a63c-3f3bd6173ea1-32_1499_1459_347_145}\\
( $e^{\prime}>e$ ) implies\\
$\left.F\left(\cdot, e^{\prime}\right)\right\rangle_{\text {FOSD }} F(\cdot, e)$\\
$(\forall y) 1-F\left(y, e^{\prime}\right)>1-F(y, e)$

\section*{Principal's problem}
$$
\max _{t(y), e} \int U_{p}(y-t(y)) f(y, e) d y
$$

s.t.


\begin{equation*}
\int U_{a}(t(y)) f(y, e) d y-\Psi(e) \geq \bar{U}_{a} \tag{IR}
\end{equation*}


(IC) $\quad e \in \arg \max _{e^{\prime} \in A}\left\{\int U_{a}(t(y)) f\left(y, e^{\prime}\right) d y-\Psi\left(e^{\prime}\right)\right\}$

The IC constraint can be replaced by the first and second order conditions of the implicit agent's maximization problem:


\begin{align*}
& \int U_{a}(t(y)) \frac{\partial f(y, e)}{\partial e} d y=\Psi^{\prime}(e)  \tag{3}\\
& \int U_{a}(t(y)) \frac{\partial^{2} f(y, e)}{\partial e^{2}} d y \leq \Psi^{\prime \prime}(e) \tag{4}
\end{align*}


\section*{First order approach}
\begin{itemize}
  \item Some authors ignore SOC (Equation 4) and use only FOC (Equation 3), for instance Holmström (1979), Shavell (1979).
  \item The first order approach is not legitimate; it may identify a minimum instead of a maximum, or the equivalent of an inflection point.
  \item We will first follow the first order approach, then illustrate why it may lead to incorrect results, and finally give conditions under which it is valid.
\end{itemize}

\section*{FOA. Program}
Replacing the IC constraint in the seller's problem by the first-order condition of the agent's problem one obtains a modified program,

$$
\begin{aligned}
& \max _{t(y), e} \int U_{p}(y-t(y)) f(y, e) d y \\
& \text { s.t. } \\
& \quad \int U_{a}(t(y)) f(y, e) d y-\Psi(e) \geq \bar{U}_{a} \\
& \quad \int U_{a}(t(y)) \frac{\partial f(y, e)}{\partial e} d y=\Psi^{\prime}(e)
\end{aligned}
$$

\section*{FOA. Lagrangean}
$$
\begin{aligned}
\mathcal{L}(t, e, \lambda, \mu)= & \int\left\{U_{p}(y-t(y)) f(y, e)\right. \\
& +\lambda\left[U_{a}(t(y)) f(y, e)-\Psi(e)-\bar{U}_{a}\right] \\
& \left.+\mu\left[U_{a}(t(y)) f_{e}(y, e)-\Psi^{\prime}(e)\right]\right\} d y
\end{aligned}
$$

where $\lambda$ and $\mu$ are the multipliers of the IR and IC constraints respectively.

FOA. First order condition, $t(\cdot)$

$$
\begin{aligned}
& -U_{p}^{\prime}(y-t(y)) f(y, e)+\lambda U_{a}^{\prime}(t(y)) f(y, e) \\
& \quad+\mu U_{a}^{\prime}(t(y)) f_{e}(y, e)=0
\end{aligned}
$$

Dividing by $U_{a}^{\prime}(t(y)) f(y, e)$ yields the necessary condition


\begin{equation*}
\frac{U_{p}^{\prime}(y-t(y))}{U_{a}^{\prime}(t(y))}=\lambda+\mu \frac{f_{e}(y, e)}{f(y, e)} . \tag{5}
\end{equation*}


\section*{Observation 1}
\begin{itemize}
  \item Grossman and Hart (1983) propose a two-step method for solving PA problems in the parametrized distribution approach.
\end{itemize}

\begin{enumerate}
  \item For each $e$, find $t(y)$ that minimizes expected $t(y)$ subject to agent's IR and IC.
  \item Find $e$ that maximizes expected output minus the expected cost.
\end{enumerate}

\section*{Observation 2}
\begin{itemize}
  \item If the "surrogate" IC constraint is not binding, $\mu=0$. Then Equation 5 becomes
\end{itemize}

$$
\frac{U_{p}^{\prime}(y-t(y))}{U_{a}^{\prime}(t(y))}=\lambda
$$

This is a necessary and sufficient condition for optimal risk sharing (Borch (1962)).

\begin{itemize}
  \item Intuition. The multiplier $\lambda$ is the utility cost to the principal of providing more utility to the agent.\\
\includegraphics[max width=\textwidth, alt={}, center]{3d1c6c7a-b07a-4a39-a63c-3f3bd6173ea1-41_1412_1418_184_150}
\end{itemize}

$$
\lambda=\frac{U_{p}^{\prime}}{U_{a}^{\prime}}=\frac{d U_{p} / d t}{d U_{a} / d t}=\frac{d U_{p}}{d U_{a}}
$$

\section*{Observation 3}
\begin{itemize}
  \item Under the maintain assumptions, if $F_{e}(y, e) \leq 0$ with strict inequality for some values of $e$, then $\mu>0$.
  \item Proof. Differentiate the Lagrangean with respect to $e$. Make use of second order condition on the Agent's problem.
  \item There is a tradeoff between optimal risk sharing and the provision of incentives for the agent.
\end{itemize}

\section*{Observation 4}
\begin{itemize}
  \item Interpretation of Equation 5 by analogy with two action case.
  \item There are only two effort levels $e_{l}<e_{h}$.
  \item If the principal wishes to implement $e=e_{l}$, the problem is simple: $t$ must be such that IR is satisfied.
  \item If the principal wishes to implement $e=e_{h}$, in addition to IR, $t$ must satisfy IC.
\end{itemize}

$$
\begin{aligned}
& \int U_{a}(t(y)) f\left(y, e_{h}\right) d y-\Psi\left(e_{h}\right) \\
& \geq \int U_{a}(t(y)) f(y, e l) d y-\Psi(e l) \\
& \int U_{a}(t(y))\left[f\left(y, e_{h}\right)-f(y, e l)\right] d y \\
& \quad-\left[\Psi\left(e_{h}\right)-\Psi(e l)\right] \geq 0
\end{aligned}
$$

\begin{itemize}
  \item Replace IC by equation above and differentiate new Lagrangean w.r.t.t.
\end{itemize}

Assuming principal wants $e=e_{h}$, then

$$
\begin{aligned}
\mathcal{L}= & \int\left\{U_{p} f\left(y, e_{h}\right)+\lambda\left[U_{a} f\left(y, e_{h}\right)-\Psi\left(e_{h}\right)-\bar{U}_{a}\right]\right. \\
& \left.+\mu \int U_{a}(t(y))\left[f\left(y, e_{h}\right)-f\left(y, e_{l}\right)\right]\right\} d y \\
& -\mu\left[\Psi\left(e_{h}\right)-\Psi\left(e_{l}\right)\right]
\end{aligned}
$$

FOC w.r.t. $t$

$$
-U_{p}^{\prime} f\left(y, e_{h}\right)+\lambda U_{a}^{\prime} f\left(y, e_{h}\right)+\mu U_{a}^{\prime}\left[f\left(y, e_{h}\right)-f\left(y, e_{l}\right)\right]=0
$$

from which one obtains the necessary condition

$$
\frac{U_{p}^{\prime}(y-t(y))}{U_{a}^{\prime}(t(y))}=\lambda+\mu\left[1-\frac{f\left(y, e_{l}\right)}{f\left(y, e_{h}\right)}\right]
$$

If principal is risk neutral, $U_{p}^{\prime}=1$. Then

$$
\frac{1}{U_{a}^{\prime}(t(y))}=\lambda+\mu\left[1-\frac{f\left(y, e_{l}\right)}{f\left(y, e_{h}\right)}\right]
$$

\begin{itemize}
  \item The equation "says" that $t(y)$ varies with $y$.
  \item When $\frac{f\left(y, e_{l}\right)}{f\left(y, e_{h}\right)}<1$ then $e_{h}$ is more likely than $e_{l}$.
\end{itemize}

$$
\frac{1}{U_{a}^{\prime}(t(y))}=\lambda+\mu\left[1-\frac{f\left(y, e_{l}\right)}{f\left(y, e_{h}\right)}\right]
$$

\begin{itemize}
  \item Intuition. Suppose higher $y\left\}\right.$ that $e_{h}$ is more likely than $e_{l}$, then $\frac{f\left(y, e_{l}\right)}{f\left(y, e_{h}\right)}$ becomes smaller.
  \item This means RHS is larger.
  \item In turn this must imply that $U_{a}^{\prime}$ is smaller.
  \item Smaller $U_{a}^{\prime}$ means larger $t(y)$ because $U_{a}$ is concave.
\end{itemize}

\section*{Observation 5}
\begin{itemize}
  \item $\frac{f\left(y, e_{l}\right)}{f\left(y, e_{h}\right)}$ is the likelihood ratio; it indicates the likelihood that the agent chooses $e_{l}$ relative to $e_{h}$.
  \item Agent is punished when outcomes move beliefs about $e_{h}$ down and rewarded in opposite case.
  \item Agency problem is not an inference problem in the statistical sense. The principal gains no new information about the action taken by the agent from observing $y$ because, by contract design, the principal already knows the action implemented.
  \item The optimal contract reflects the inference principle.
\end{itemize}

\section*{Observation 6}
\begin{center}
\begin{tabular}{llll}
$f(y, e)$ & $y_{l}$ & $y_{m}$ & $y_{h}$ \\
\hline
$e_{h}$ & 0 & 0.5 & 0.5 \\
\hline
$e_{l}$ & 0.5 & 0.1 & 0.4 \\
\hline
$\frac{f\left(y, e_{l}\right)}{f\left(y, e_{h}\right)}$ & $\infty$ & 0.20 & 0.8 \\
\hline
\end{tabular}
\end{center}

\includegraphics[max width=\textwidth, alt={}, center]{3d1c6c7a-b07a-4a39-a63c-3f3bd6173ea1-49_866_797_1045_153}\\
$t(y)$ depends on likelihood ratio. Even under FOSD optimal contract need not be monotone.

Optimal contract: $t(y l)<t(y h)<t(y m)$

\section*{Monotone optimal contract}
\begin{itemize}
  \item FOC condition for optimality is
\end{itemize}

$$
\frac{1}{U_{a}^{\prime}(t(y))}=\lambda+\mu\left[1-\frac{f\left(y, e_{l}\right)}{f\left(y, e_{h}\right)}\right]
$$

\begin{itemize}
  \item $t(y)$ is monotone nondecreasing in $y$ if and only if $\frac{d\left[\frac{f\left(y, e_{l}\right)}{f\left(y, e_{h}\right)}\right]}{d y} \leq 0$.
\end{itemize}

\section*{Monotone and linear contracts}
\begin{itemize}
  \item Monotone contracts require MLRP, a strong assumption.
  \item If agent can "reduce" output at no cost/effort, then transfer function ought to be monotone; otherwise the agent would artificially reduce output when necessary.
\end{itemize}

In example: reduce $y_{h}$ to $y_{m}$ when the outcome is $y_{h}$.

\begin{itemize}
  \item Linear contracts
  \item 
1.


  \item 2 .
\end{itemize}

\section*{Value of Information: Statistics preliminaries}
Definition 5 Let $S \subset \mathbb{R}^{n}$ be a sample space. A model is a family of probability distributions $\left\{P_{\theta}: \theta \in \Theta\right\}$ where $(\forall \theta \in \Theta) P_{\theta}$ is a probability distribution on $S$ and $\Theta$ is a set, called the parameter set.

A model is regular if ( $\forall \theta \in \Theta$ ) either

\begin{enumerate}
  \item $P_{\theta}$ is $\mathbb{C}^{0}$ and has density $p(x, \theta)$.
  \item $P_{\theta}$ is discrete with frequency $p(x, \theta)$ and $\exists\left\{x_{1}, x_{2}, \ldots\right\} \subseteq S$ such that $\sum_{i=1}^{\infty} p\left(x_{i}, \theta\right)=1$.
\end{enumerate}

Note: (1) requires that every $P_{\theta}$ be absolutely continuous.

\section*{Statistic preliminaries}
Definition 6 A statistic in the model is any function $T: S \rightarrow W$ where $W$ is a space.

\begin{itemize}
  \item $W$ is often some Euclidean space or even $\mathbb{R}$.
\end{itemize}

Definition 7 A statistic $T$ is a reduction of data if ( $\exists x, x^{\prime} \in S, x \neq x^{\prime}$ ) and $T(x)=T\left(x^{\prime}\right)$.

\begin{itemize}
  \item Any statistic that is a reduction of data causes a loss of information.
  \item A statistic $T$ is sufficient for a parameter of interest, if observing the realization $T(x)=t$, does not imply a loss of information about the parameter of interest.
\end{itemize}

\section*{Statistic preliminaries}
Definition 8 The statistics $T: S \rightarrow W$ is sufficient for $\theta$ if $\left(\exists Q: \mathbb{R}^{n}\right. \times W \rightarrow[0,1])(\forall \theta \in \Theta)(\forall t \in W)\left[P_{\theta}(x \mid T(x)=t)=Q(x, t)\right]$.

\begin{itemize}
  \item The probability of $x$ given $T(x)=t$, does not depend on $\theta$. Equivalently, once $t$ is observed, knowing the sample realization $x$ (such $T(x)=t$ ) does not contain additional information about $\theta$ (or $P_{\theta}$ ) provided the model is valid.
  \item Determining directly whether a statistic is sufficient for a parameter is often difficult because conditional probabilities must be computed.
\end{itemize}

\section*{Statistic preliminaries}
Theorem 4 (Factorization theorem) $T$ is sufficient for $\theta$ if and only if ( $\exists g : W \times \Theta \rightarrow \mathbb{R})(\exists h: S \rightarrow \mathbb{R})$ such that $(\forall x \in S)(\forall \theta \in \Theta)$

$$
p(x, \theta)=g(T(x), \theta) h(x)
$$

\begin{itemize}
  \item The Factorization theorem holds in more general environments than the regular models presented here. For a general treatment, see Lehmann (1997).
\end{itemize}

\section*{Value of information}
\begin{itemize}
  \item The optimal contract exploits the correlation between $y$ and $e$. Suppose there is another variable $x$ that is correlated with $e$.
  \item Would a transfer $t(x, y)$ dominate the optimal $t(y)$ ? Yes, whenever $x$ gives information about $e$ that is not contained in $y$.
\end{itemize}

This results is due to Holmström (1979).

\section*{Value of information in principal agent model}
Theorem 5 Let $t(y)$ be an optimal transfer function for which the agent's choice of action is unique and interior in $A$. Then, there exist $t(x, y)$ (where $x$ is another observable variable) which strictly pareto dominates $t(y)$ if and only if $x$ is informative about $e$.

Definition $9 x$ is informative about $e$ whenever $y$ is not a sufficient statistic for $(x, y)$ with respect to $e$.

The result is due to Holmström (1979); the paper contains a useful discussion and the proof of the theorem.

\section*{Mirrlees and the FOA}
\begin{itemize}
  \item See Mirrlees (1999). This paper was written in 1975. It was broadly read as a working paper.
  \item See Rogerson (1985).
  \item See Jewitt (1988).
  \item See Conlon (2009).
  \item See Chade and Swinkels (2019).
\end{itemize}

\section*{References}
Bolton, Patrick, and Dewatripont. 2005. Contract Theory. MIT Press.\\
Borch, Karl. 1962. "Equilibrium in a Reinsurance Market." Econometrica 30 (3): 424-44.\\
\href{https://doi.org/10.2307/1909887}{https://doi.org/10.2307/1909887}.\\
Chade, Hector, and Jeroen Swinkels. 2019. "The No-Upward-Crossing Condition, Comparative Statics, and the Moral Hazard Problem." Unpublished manuscript.\\
Conlon, J. 2009. "Two New Conditions Supporting the First-Order Approach to Multisignal PrincipalAgent Problems." Econometrica 77: 249-78.\\
Demski, Joel S., and David M. Kreps. 1982. "Models in Managerial Accounting." Journal of Accounting Research 20: 117-48.\\
Diamond, Douglas W. 1984. "Financial Intermediation and Delegated Monitoring." The Review of Economic Studies 51: 393-414. \href{https://doi.org/https://doi.org/10.2307/2297430}{https://doi.org/https://doi.org/10.2307/2297430}.\\
Grossman, Sanford, and Oliver Hart. 1983. "An Analysis of the Principal-Agent Problem." Econometrica 51: 7-45.\\
Hart, Oliver, and Bengt Holmström. 1987. "The Theory of Contracts." In Advances in Economic Theory, Fifth World Congress, edited by Truman Bewley. Cambridge University Press.\\
Holmström, Bengt. 1979. "Moral Hazard and Observability." Bell Journal of Economics 10: 74-91. Jensent, Michael C., and William H. Meckling. 1976. Jounral of Finance 3: 305-60.

Jewitt, Ian. 1988. "Justifying the First-Order Approach to Principal-Agent Problems." Econometrica, 1177-90.\\
Lehmann, E. L. 1997. Testing Statistical Hypotheses. 2nd ed. Springer Texts in Statistics. Springer. \href{https://doi.org/10.1007/978-1-4757-1923-9}{https://doi.org/10.1007/978-1-4757-1923-9}.\\
Marshall, Albert W., Ingram Olkin, and Barry C. Arnold. 2010. Inequalities: Theory of Majorization and Its Applications. 2nd ed. Springer. \href{https://doi.org/10.1007/978-0-387-68276-1}{https://doi.org/10.1007/978-0-387-68276-1}.\\
Mas-Colell, Andreu, Michael Whinston, and Jerry Green. 1995. Microeconomic Theory. Oxford University Press.\\
Milgrom, Paul R. 1981. "Good News and Bad News: Representation Theorems and Applications." The Bell Journal of Economics 12 (2): 380-91.\\
Mirrlees, James A. 1999. "The Theory of Moral Hazard and Unobservable Behaviour: Part i." The Review of Economic Studies 66 (1): 3-21.\\
Rogerson, W. 1985. "The First-Order Approach to Principal-Agents Problems." Econometrica 53: 135767.

Rothschild, Michael, and Joseph Stiglitz. 1976. "Equilibrium in Competitive Insurance Markets: An Essay on the Economics of Imperfect Information." The Quarterly Journal of Economics 90 (4): 629-49. \href{http://www.jstor.org/stable/1885326}{http://www.jstor.org/stable/1885326}.\\
Royden, Halsey Lawrence. 1968. Real Analysis. 2nd ed. McMillan Publishing Co. Shavell, Steven. 1979. "On Moral Hazard and Insurance*." The Quarterly Journal of Economics 93 (4): 541-62. \href{https://doi.org/10.2307/1884469}{https://doi.org/10.2307/1884469}.


\end{document}